\section{Discussion and conclusions}
\subsection{Ethical implications}
The stereotype task highlights the importance of defining what a fairness detector
is intended to recognize. A system that only detects explicit generalizations may
miss implicit occupational or religious associations. Conversely, treating every
positive portrayal as a stereotype may suppress useful counter-stereotypical
representation. These classification results support evaluating the detector's
categories and reviewing concrete errors; they do not establish that the model
treats all demographic groups equitably in hiring, education, or other decisions.

Robustness matters when users write imperfect English or when an attacker edits
incoming material. A change in wording can alter a downstream decision even if the
reference label is unchanged. However, requiring invariance after a real change
in meaning would also be unsafe. Several selected attacks replace concepts or
question types, so a changed answer may be a defensible semantic judgment rather
than a successful manipulation. The reported ASR measures reference-label failures
in the supplied dataset, and should not be interpreted as a verified rate of
meaning-preserving attacks against a deployed application.

Moral action classification can reveal whether the model identifies the action,
actor, and context relevant to a judgment. It remains a test of agreement with a
dataset's binary judgments, rather than a certificate of moral competence. A refusal
to judge or an overly broad condemnation can also fail a benchmark without providing
evidence that the system encourages harmful actions. Explanations are useful for
inspection but are generated text, not verified accounts of the model's internal
decision process. Automated ethical decisions would require contextual review and
accountability beyond these scores.

\subsection{Limitations and future work}
Only one quantized model, one runtime configuration, and three English subsets were
evaluated. The study omits other TrustLLM dimensions and other tasks within the
selected dimensions, as well as multilingual, multimodal, and open-ended behavior.
The JSON constraint and explanation request change the original evaluation protocol;
different prompting, thinking settings, or quantization may change results. A single
run does not estimate variation across seeds or deployments. Prompt-level overlap
in AdvGLUE and unequal fairness-topic coverage further constrain generalization.
The public datasets may have appeared in training; no contamination audit was done.
No confidence intervals are claimed from the correlated paired attack sample.

Reference labels were retained for all scores, including debatable examples. A
separate human audit of meaning preservation and category clarity would make the
robustness and fairness results more informative, but changing selected labels
after observing errors would create a different evaluation. Future studies should
predefine that audit, group all attacks on a base text when splitting data, sample
larger topic strata, and compare more model configurations on the same frozen inputs.

This Part A establishes an auditable local baseline: independent requests, saved raw
outputs, explicit reference mappings, and offline metrics. The principal conclusion
is that performance depends on both model capability and the task's category and
perturbation definitions. Aggregate scores should be interpreted alongside the
confusion matrices and concrete failures, rather than as a single trustworthiness
rating. No Part B mitigation was applied in this baseline evaluation.

\subsection*{Use of generative AI}
Generative AI assistance (Codex) was used to discuss the evaluation design, implement
and check the Python scripts, inspect outputs, and prepare this report draft. The
benchmark responses were generated by the local Gemma model under the recorded
configuration; its generated explanations were not used as an external judge.
Scores were computed programmatically against saved reference labels. The draft
includes AI-assisted interpretations of selected failures.
